\chapter{Geometric Lagrangian And Hamiltonian Mechanics}
\label{chap:lag-ham}

\section{Mechanical Data}
\label{sec:lag-ham-mechanical-data}

The first question in any mechanics problem is not ``what is the force?'' but
``what is the space of possible configurations?'' Once that space is named, the
rest of the theory becomes coordinate-independent.

\begin{definition}[Finite-dimensional mechanical system]
A regular finite-dimensional Lagrangian mechanical system consists of:
\begin{enumerate}
\item a smooth \(n\)-dimensional configuration manifold \(Q\);
\item a time interval \(I=[t_0,t_1]\);
\item a Lagrangian \(L:TQ\times I\to \R\);
\item a class of admissible curves \(q:I\to Q\), possibly satisfying constraints.
\end{enumerate}
The action of an admissible curve is
\[
  S[q]=\int_{t_0}^{t_1} L(q(t),\dot q(t),t)\,\dd t.
\]
\end{definition}

\begin{assumption}[Coordinate patch for derivations]
Derivations in this chapter are written in one coordinate chart
\((q^1,\ldots,q^n)\). The equations obtained are tensorial or intrinsic when
the objects used to build them are intrinsic. Coordinates are a calculation
device, not an additional physical structure.
\end{assumption}

\begin{table}[htbp]
\centering
\caption{Core notation for Chapter~\ref{chap:lag-ham}. Global conventions and
overloaded symbols are listed in Chapter~\ref{chap:notation}.}
\label{tab:lag-ham-notation}
\small
\begin{tabular}{p{0.18\textwidth}p{0.45\textwidth}p{0.28\textwidth}}
\toprule
Symbol & Definition & Interpretation or first use\\
\midrule
\(Q\) & smooth configuration manifold & space of positions\\
\(I=[t_0,t_1]\) & time interval & domain of a path\\
\(q:I\to Q\) & configuration path & physical motion candidate\\
\(q^i(t)\) & local coordinates of \(q(t)\) & generalized positions\\
\(\dot q^i(t)\) & \(\dd q^i/\dd t\) & generalized velocities\\
\(\eta^i(t)\) & \(\delta q^i(t)\) & infinitesimal virtual displacement\\
\(L:TQ\times I\to\R\) & Lagrangian & input to Hamilton's principle\\
\(S[q]\) & \(\int L\,\dd t\) & action functional\\
\(p_i\) & \(\partial L/\partial \dot q^i\) & canonical momentum\\
\(E_L\) & \(p_i\dot q^i-L\) & Lagrangian energy\\
\(g_{ij}\) & mass metric & natural Lagrangian \(T-V\)\\
\(\Gamma^\ell{}_{jk}\) & Christoffel symbols of \(g\) & inertial terms in coordinates\\
\(\Phi^a(q,t)\) & holonomic constraints & admissible submanifold\\
\(\lambda_a\) & constraint multipliers & reaction forces\\
\(G,\mathfrak g\) & symmetry group and Lie algebra & Noether theorem\\
\(\xi_Q\) & infinitesimal generator & vector field on \(Q\)\\
\(J_\xi\) & \(p_i\xi_Q^i\) & Noether charge\\
\(r,\rho,\theta\) & Kepler vector, radius, polar angle & central-force example\\
\(\ell\) & \(\mu\rho^2\dot\theta\) & angular momentum magnitude\\
\(R\) & attitude matrix & rigid-body example\\
\(\Omega,M\) & body angular velocity and body angular momentum & \(M=I\Omega\)\\
\(H\) & \(p_i\dot q^i-L\) after Legendre transform & Hamiltonian\\
\(\Theta\) & \(p_i\,\dd q^i\) & canonical one-form\\
\(\omega\) & \(-\dd\Theta=\dd q^i\wedge\dd p_i\) & symplectic form\\
\(X_H\) & \(\iota_{X_H}\omega=\dd H\) & Hamiltonian vector field\\
\(\Phi_H^t\) & Hamiltonian flow map & phase-space evolution\\
\(\Psi_{\Delta t}\) & numerical one-step map & integrator output\\
\bottomrule
\end{tabular}
\end{table}

\paragraph{Roadmap.}
The chapter has three jobs. First it derives the Euler-Lagrange equations from
Hamilton's principle, including boundary terms, constraints, and Noether
charges. Second it treats Kepler motion and the rigid body as the first serious
examples of the formalism. Third it passes through the Legendre transform to
Hamiltonian mechanics, where phase space, symplectic form, Poisson brackets,
and canonical transformations become the basic language. The point is not to
replace Newton's law by harder notation; it is to make clear which parts of
mechanics are coordinate choices and which parts are geometric structure.

The chapter should be read as a change in viewpoint. In the Lagrangian picture,
the primitive object is a path in configuration space and the variational
principle selects the physical path. In the Hamiltonian picture, the primitive
object is a point in phase space and the symplectic form turns the Hamiltonian
function into a vector field. The two pictures are equivalent only under the
regularity assumption stated in Chapter~\ref{chap:notation}; this is why the
Legendre transform is a theorem-sized step, not a change of notation.

\paragraph{Derivation checkpoints.}
Whenever this chapter differentiates an action, check three things before
following the algebra: which variables are varied, which endpoint conditions
are imposed, and which boundary terms survive. Whenever it changes from
\((q,\dot q)\) to \((q,p)\), check the velocity Hessian. Whenever it invokes a
symmetry, check whether the symmetry acts on configuration space, phase space,
or both. These checks are small, but they prevent the most common mistakes:
using an inadmissible virtual displacement, assuming a singular Legendre map is
invertible, or calling a quantity conserved before verifying the symmetry
hypothesis.

\section{Variations And The Euler-Lagrange Equations}
\label{sec:euler-lagrange}

Let \(q_\eps(t)\) be a smooth one-parameter family of curves with
\(q_0(t)=q(t)\). The variation vector field along \(q\) is
\[
  \eta(t)=\left.\frac{\partial q_\eps(t)}{\partial \eps}\right|_{\eps=0},
  \qquad
  \eta^i(t)=\left.\frac{\partial q_\eps^i(t)}{\partial \eps}\right|_{\eps=0}.
\]
For fixed-endpoint variations we impose
\[
  \eta(t_0)=0,\qquad \eta(t_1)=0.
\]
The variational calculation below has one essential idea. The path \(q(t)\) is
a critical point of the action when every infinitesimal deformation of the path
changes the action only at second order. Differentiating the action produces a
term proportional to \(\dot\eta\); integration by parts moves the derivative
onto the momentum, and the endpoint condition kills the boundary term. What is
left must vanish for every interior virtual displacement \(\eta\), which forces
the Euler-Lagrange equations pointwise in time.

\begin{figure}[htbp]
\centering
\begin{tikzpicture}[scale=1.0, >=Latex]
  \draw[->] (-0.3,0) -- (7.2,0) node[right] {\(t\)};
  \draw[->] (0,-0.2) -- (0,3.0) node[above] {configuration coordinate};
  \draw[densely dotted] (0.5,0) -- (0.5,2.75);
  \draw[densely dotted] (6.5,0) -- (6.5,2.75);
  \draw[thick, blue] plot[smooth] coordinates {(0.5,0.7) (1.8,1.25) (3.2,1.0) (4.7,2.0) (6.5,2.35)};
  \draw[thick, red, dashed] plot[smooth] coordinates {(0.5,0.7) (1.8,1.55) (3.2,1.25) (4.7,1.75) (6.5,2.35)};
  \draw[->, red] (3.2,1.0) -- (3.2,1.25) node[midway,right] {\(\eta(t)\)};
  \draw[->, red!70!black] (4.7,2.0) -- (4.7,1.75);
  \fill (0.5,0.7) circle (2pt) node[below right] {\(q(t_0)\)};
  \fill (6.5,2.35) circle (2pt) node[above left] {\(q(t_1)\)};
  \node[below] at (0.5,0) {\(t_0\)};
  \node[below] at (6.5,0) {\(t_1\)};
  \node[blue] at (2.3,0.6) {\(q(t)\)};
  \node[red] at (4.25,1.35) {\(q_\eps(t)=q(t)+\eps\eta(t)+O(\eps^2)\)};
\end{tikzpicture}
\caption{A fixed-endpoint variation changes the path in the interior while
holding the initial and final configurations fixed. The infinitesimal change is
the vector field \(\eta(t)=\delta q(t)\) along the path.}
\label{fig:fixed-endpoint-variation}
\end{figure}

\begin{derivation}[Euler-Lagrange equations]
The first variation of the action is
\[
\begin{aligned}
  \delta S
  &=
  \left.\frac{\dd}{\dd\eps}\right|_{\eps=0}
  \int_{t_0}^{t_1}L(q_\eps,\dot q_\eps,t)\,\dd t\\
  &=
  \int_{t_0}^{t_1}
  \left(
    \frac{\partial L}{\partial q^i}\eta^i
    +
    \frac{\partial L}{\partial \dot q^i}\dot\eta^i
  \right)\dd t.
\end{aligned}
\]
Integrate the second term by parts:
\[
  \int_{t_0}^{t_1}
  \frac{\partial L}{\partial \dot q^i}\dot\eta^i\,\dd t
  =
  \left[
    \frac{\partial L}{\partial \dot q^i}\eta^i
  \right]_{t_0}^{t_1}
  -
  \int_{t_0}^{t_1}
  \frac{\dd}{\dd t}\left(\frac{\partial L}{\partial \dot q^i}\right)\eta^i
  \,\dd t.
\]
The boundary term vanishes because \(\eta(t_0)=\eta(t_1)=0\). Therefore
\[
  \delta S
  =
  \int_{t_0}^{t_1}
  \left[
    \frac{\partial L}{\partial q^i}
    -
    \frac{\dd}{\dd t}\left(\frac{\partial L}{\partial \dot q^i}\right)
  \right]\eta^i\,\dd t.
\]
The fundamental lemma of the calculus of variations says that if this integral
vanishes for all compactly supported variations \(\eta\), then the coefficient
of every \(\eta^i\) vanishes. Hence
\[
  \boxed{
  \frac{\dd}{\dd t}\left(\frac{\partial L}{\partial \dot q^i}\right)
  -
  \frac{\partial L}{\partial q^i}=0
  }.
\]
\end{derivation}

The boxed equation is local in time because the variations are arbitrary in the
interior of the interval. This is the conceptual content of the fundamental
lemma: a single global statement, \(\delta S=0\), becomes a pointwise
differential equation only because the virtual displacement \(\eta(t)\) can be
chosen independently near each time. Boundary conditions and endpoint terms are
therefore not afterthoughts; they decide which variations are admissible and
which physical quantities appear at the endpoints.

\section{Boundary Terms And Forces}
\label{sec:boundary-forces}

The integration-by-parts calculation also explains momenta. Without imposing
fixed endpoints, the boundary contribution is
\[
  \left[p_i\eta^i\right]_{t_0}^{t_1},
  \qquad
  p_i=\frac{\partial L}{\partial \dot q^i}.
\]
Thus \(p_i\) is the coefficient of endpoint displacement in the first variation
of the action. This is a sharper definition than ``mass times velocity'', which
is only true in special coordinates for special Lagrangians.

If nonconservative generalized forces \(Q_i(q,\dot q,t)\) are present, the
Lagrange-d'Alembert principle is
\[
  \delta S+\int_{t_0}^{t_1}Q_i\eta^i\,\dd t=0,
\]
which gives
\[
  \frac{\dd}{\dd t}\left(\frac{\partial L}{\partial \dot q^i}\right)
  -
  \frac{\partial L}{\partial q^i}
  =
  Q_i.
\]
The present course focuses mainly on conservative and Hamiltonian systems, but
this formula is useful when comparing with fluids and continua.

\section{Natural Lagrangians And The Metric Form}
\label{sec:natural-lagrangians}

Many mechanical systems have a Lagrangian of the form
\[
  L(q,\dot q)=T(q,\dot q)-V(q)
  =
  \frac{1}{2}g_{ij}(q)\dot q^i\dot q^j - V(q),
\]
where \(g_{ij}\) is a positive-definite mass metric on \(Q\), and \(V\) is
potential energy.

\begin{derivation}[Geodesic-plus-force equation]
For the natural Lagrangian,
\[
  \frac{\partial L}{\partial \dot q^i}=g_{ij}\dot q^j,
\]
so
\[
  \frac{\dd}{\dd t}\frac{\partial L}{\partial \dot q^i}
  =
  g_{ij}\ddot q^j
  +
  \frac{\partial g_{ij}}{\partial q^k}\dot q^k\dot q^j.
\]
Also
\[
  \frac{\partial L}{\partial q^i}
  =
  \frac{1}{2}\frac{\partial g_{jk}}{\partial q^i}\dot q^j\dot q^k
  -
  \frac{\partial V}{\partial q^i}.
\]
The Euler-Lagrange equation becomes
\[
  g_{ij}\ddot q^j
  +
  \left(
    \frac{\partial g_{ij}}{\partial q^k}
    -
    \frac{1}{2}\frac{\partial g_{jk}}{\partial q^i}
  \right)\dot q^k\dot q^j
  =
  -\frac{\partial V}{\partial q^i}.
\]
After symmetrizing the velocity terms and raising an index with \(g^{\ell i}\),
\[
  \boxed{
  \ddot q^\ell+\Gamma^\ell_{jk}\dot q^j\dot q^k
  =
  -g^{\ell i}\frac{\partial V}{\partial q^i}
  },
\]
where
\[
  \Gamma^\ell_{jk}
  =
  \frac{1}{2}g^{\ell i}
  \left(
    \frac{\partial g_{ij}}{\partial q^k}
    +
    \frac{\partial g_{ik}}{\partial q^j}
    -
    \frac{\partial g_{jk}}{\partial q^i}
  \right)
\]
are the Christoffel symbols of the mass metric.
\end{derivation}

Thus even ``force-free'' motion on a curved configuration space need not look
like \(\ddot q=0\) in coordinates. The Christoffel terms are inertial terms
created by the geometry of \(Q\) and the chosen coordinates.

\section{Holonomic Constraints}
\label{sec:holonomic-constraints}

A holonomic constraint is a condition
\[
  \Phi^a(q,t)=0,\qquad a=1,\ldots,k,
\]
that cuts out an admissible submanifold of configuration space. If one works in
the ambient coordinates \(q^i\), the constrained equations take the multiplier
form
\[
  \frac{\dd}{\dd t}\frac{\partial L}{\partial \dot q^i}
  -
  \frac{\partial L}{\partial q^i}
  =
  \lambda_a\frac{\partial \Phi^a}{\partial q^i}.
\]
The multipliers \(\lambda_a\) are not arbitrary forces. They are whatever
constraint reactions are needed to keep the motion on \(\Phi^a=0\). If instead
one parametrizes the constraint surface directly, the \(\lambda_a\) disappear
and the same physics is expressed with fewer coordinates.

\section{Symmetry And Noether's Theorem}
\label{sec:noether}

Noether's theorem is the first place where the variational viewpoint pays a
large dividend. A symmetry is not merely a visual invariance of a picture; it
is a transformation of the configuration space that leaves the Lagrangian, and
therefore the action, unchanged. The infinitesimal version of that
transformation is a vector field \(\xi_Q\) along which the action has zero
first variation. The conserved quantity is the momentum paired with that vector
field.

\begin{definition}[Infinitesimal generator]
Let a Lie group \(G\) act on \(Q\). For \(\xi\in\mathfrak g\), the infinitesimal
generator is the vector field
\[
  \xi_Q(q)=\left.\frac{\dd}{\dd s}\right|_{s=0}\exp(s\xi)\cdot q.
\]
In coordinates, write \(\xi_Q=\xi_Q^i(q)\partial/\partial q^i\).
\end{definition}

\begin{assumption}[Lagrangian symmetry]
The action of \(G\) is a symmetry if
\[
  L(g\cdot q,\,T_qg\cdot \dot q,\,t)=L(q,\dot q,t)
\]
for all \(g\in G\), \(q\in Q\), and \(\dot q\in T_qQ\).
\end{assumption}

\begin{proposition}[Noether theorem]
For every \(\xi\in\mathfrak g\), the quantity
\[
  J_\xi(q,\dot q)
  =
  \frac{\partial L}{\partial \dot q^i}\xi_Q^i(q)
  =
  p_i\xi_Q^i(q)
\]
is conserved along solutions of the Euler-Lagrange equations.
\end{proposition}

\begin{proof}
Differentiate the symmetry condition along \(g=\exp(s\xi)\) at \(s=0\):
\[
  0=
  \frac{\partial L}{\partial q^i}\xi_Q^i
  +
  \frac{\partial L}{\partial \dot q^i}
  \frac{\dd}{\dd t}\xi_Q^i(q(t)).
\]
Along an Euler-Lagrange solution,
\[
  \frac{\partial L}{\partial q^i}
  =
  \frac{\dd}{\dd t}\frac{\partial L}{\partial \dot q^i}.
\]
Substitute this into the previous equation:
\[
  0=
  \frac{\dd}{\dd t}\left(\frac{\partial L}{\partial \dot q^i}\right)\xi_Q^i
  +
  \frac{\partial L}{\partial \dot q^i}
  \frac{\dd}{\dd t}\xi_Q^i
  =
  \frac{\dd}{\dd t}
  \left(
    \frac{\partial L}{\partial \dot q^i}\xi_Q^i
  \right).
\]
\end{proof}

\begin{example}[Translation and rotation]
For a particle in \(\R^3\) with \(L=\frac12m\norm{\dot x}^2-V(x)\), translation
symmetry in direction \(a\) gives \(J_a=m\dot x\cdot a\). If \(V\) is invariant
under all translations, linear momentum is conserved. Rotational symmetry about
axis \(a\) gives
\[
  J_a=(x\times m\dot x)\cdot a,
\]
so angular momentum is conserved when \(V\) is rotationally invariant.
\end{example}

\section{Basic Example: The Kepler Problem}
\label{sec:kepler-lagrangian}

The Kepler problem should appear immediately after the Lagrangian formalism,
because it is the simplest serious payoff of cyclic coordinates, Noether's
theorem, and reduction by symmetry.

\begin{assumption}[Kepler model]
We study a particle of reduced mass \(\mu\) moving in \(\R^3\setminus\{0\}\)
under an attractive inverse-square central potential. The Lagrangian is
\[
  L(r,\dot r)
  =
  \frac12\mu\norm{\dot r}^2+\frac{k}{\norm r},
  \qquad
  k=Gm_1m_2>0.
\]
The coordinate \(r=x_1-x_2\) is the relative separation of the two bodies.
\end{assumption}

Rotational invariance gives conservation of angular momentum
\[
  J=r\times p=\mu r\times\dot r.
\]
Since \(J\cdot r=0\) and \(J\cdot \dot r=0\), the motion stays in the plane
perpendicular to \(J\). In polar coordinates \((\rho,\theta)\) in that plane,
\[
  L=\frac12\mu(\dot\rho^2+\rho^2\dot\theta^2)+\frac{k}{\rho}.
\]
The angle \(\theta\) is cyclic, hence
\[
  \ell=\frac{\partial L}{\partial\dot\theta}=\mu\rho^2\dot\theta
\]
is conserved. The energy is
\[
  E=\frac12\mu\dot\rho^2+\frac{\ell^2}{2\mu\rho^2}-\frac{k}{\rho}.
\]
Thus the radial motion is one-dimensional motion in the effective potential
\[
  V_{\mathrm{eff}}(\rho)=\frac{\ell^2}{2\mu\rho^2}-\frac{k}{\rho}.
\]
In this subsection \(\ell\) is the ordinary, massful angular momentum of the
relative two-body problem. In Chapter~\ref{chap:perturbation-chaos}, where
asteroid motion is written per unit asteroid mass, the same orbit equation is
often written using the specific angular momentum
\(\ell_{\mathrm{spec}}=\ell/\mu=\rho^2\dot\theta\). The detailed central-force
derivation in Section~\ref{subsec:detailed-binet} keeps an arbitrary particle
mass \(m\), making the conversion between the two conventions explicit.

\begin{figure}[htbp]
\centering
\begin{tikzpicture}[scale=1.0, >=Latex]
  \draw[->] (0,0) -- (6.5,0) node[right] {\(\rho\)};
  \draw[->] (0,-1.7) -- (0,3.1) node[above] {\(V_{\mathrm{eff}}\)};
  \fill[blue!8] (0.624,-0.7) rectangle (2.52,0);
  \draw[blue, thick, domain=0.36:6.1, samples=160]
    plot (\x,{1.1/(\x*\x)-2.2/\x});
  \draw[dashed] (0,-0.7) -- (6.2,-0.7) node[right] {bound \(E<0\)};
  \draw[dashed] (0,0.45) -- (6.2,0.45) node[right] {scattering \(E>0\)};
  \node[blue] at (4.45,-1.22) {\(V_{\mathrm{eff}}\)};
  \fill (1.0,-1.1) circle (1.4pt) node[below] {minimum};
  \fill (0.624,-0.7) circle (1.3pt);
  \fill (2.52,-0.7) circle (1.3pt);
  \draw[<->, blue!70!black] (0.624,-1.42) -- (2.52,-1.42)
    node[midway, below] {allowed radial interval};
  \node[align=center] at (2.55,2.05) {centrifugal barrier\\near \(\rho=0\)};
\end{tikzpicture}
\caption{The Kepler problem reduces to one-dimensional radial motion in an
effective potential. For a fixed negative energy, the radial coordinate is
trapped between two turning points; for positive energy there is only a closest
approach before scattering.}
\label{fig:kepler-effective-potential}
\end{figure}

\begin{derivation}[Orbit equation]
Set \(u=1/\rho\). Since \(\ell=\mu\rho^2\dot\theta\),
\[
  \frac{\dd}{\dd t}=\dot\theta\frac{\dd}{\dd\theta}
  =
  \frac{\ell}{\mu}\,u^2\frac{\dd}{\dd\theta}.
\]
The radial Euler-Lagrange equation is
\[
  \mu(\ddot\rho-\rho\dot\theta^2)=-\frac{k}{\rho^2}.
\]
In terms of \(u(\theta)\) this becomes
\[
  \frac{\dd^2u}{\dd\theta^2}+u=\frac{\mu k}{\ell^2}.
\]
Therefore
\[
  u(\theta)=\frac{\mu k}{\ell^2}\left[1+e\cos(\theta-\theta_0)\right],
\]
or
\[
  \boxed{
  \rho(\theta)=
  \frac{\ell^2/(\mu k)}{1+e\cos(\theta-\theta_0)}
  }.
\]
The parameter \(e\) is the eccentricity. The conic is an ellipse, parabola, or
hyperbola according as \(e<1\), \(e=1\), or \(e>1\).
\end{derivation}

The Kepler problem also has the Laplace-Runge-Lenz vector
\[
  A=p\times J-\mu k\frac{r}{\norm r},
\]
which is conserved and points toward periapse. Its existence is special: it is
the extra conservation law behind the absence of perihelion precession in the
ideal inverse-square problem.

\section{Basic Example: Rigid Body As A Lagrangian System}
\label{sec:rigid-body-lagrangian}

The second basic example is the rigid body. It is best introduced directly after
Lagrangian mechanics, before using it later as a reduced Hamiltonian and
integrable system.

\begin{assumption}[Free rigid body]
A material point \(X\) in the body frame is located at
\[
  x(t)=R(t)X,
  \qquad R(t)\in\SO(3),
\]
with the center of mass fixed at the origin. The body is rigid, so all dynamics
is in the attitude \(R(t)\). No external torque acts.
\end{assumption}

The body angular velocity is defined by
\[
  R^{-1}\dot R=\widehat\Omega,
  \qquad
  \widehat\Omega v=\Omega\times v.
\]
The kinetic energy is
\[
  L(R,\dot R)=T
  =
  \frac12\Omega^TI\Omega,
\]
where \(I\) is the inertia tensor in the body frame. In principal axes,
\[
  I=\diag(I_1,I_2,I_3),
  \qquad
  L=\frac12(I_1\Omega_1^2+I_2\Omega_2^2+I_3\Omega_3^2).
\]

\begin{derivation}[Euler equation from constrained variations]
A variation inside \(\SO(3)\) has the form
\[
  \delta R=R\widehat\Sigma,
\]
where \(\Sigma(t_0)=\Sigma(t_1)=0\). This implies
\[
  \delta\Omega=\dot\Sigma+\Omega\times\Sigma.
\]
Let \(M=\partial L/\partial\Omega=I\Omega\). Then
\[
\begin{aligned}
  \delta S
  &=
  \int M\cdot(\dot\Sigma+\Omega\times\Sigma)\,\dd t\\
  &=
  \int(-\dot M+M\times\Omega)\cdot\Sigma\,\dd t.
\end{aligned}
\]
Stationarity gives
\[
  \dot M=M\times\Omega,
\]
or, with the hat convention \(\widehat\Omega v=\Omega\times v\) used
throughout,
\[
  \dot M+\Omega\times M=0.
\]
In principal axes this is
\[
\begin{aligned}
  I_1\dot\Omega_1 &= (I_2-I_3)\Omega_2\Omega_3,\\
  I_2\dot\Omega_2 &= (I_3-I_1)\Omega_3\Omega_1,\\
  I_3\dot\Omega_3 &= (I_1-I_2)\Omega_1\Omega_2.
\end{aligned}
\]
\end{derivation}

This derivation is deliberately placed here: the rigid body is not first of all
an abstract Lie-Poisson system. It is a Lagrangian system on \(\SO(3)\), and the
Lie-Poisson picture is what appears after symmetry reduction.

\section{Legendre Transform And Hamiltonian}
\label{sec:legendre-hamiltonian}

The Legendre transform sends \((q,\dot q)\in TQ\) to \((q,p)\in T^*Q\) by
\[
  p_i=\frac{\partial L}{\partial \dot q^i}.
\]
For a regular Lagrangian this can be inverted to express \(\dot q^i\) as a
function of \((q,p,t)\). Concretely, regularity means that the velocity Hessian
\[
  W_{ij}(q,\dot q,t)=
  \frac{\partial^2L}{\partial\dot q^i\partial\dot q^j}
\]
is nonsingular. The inverse function theorem then says that, locally in phase
space, the equations \(p_i=\partial L/\partial\dot q^i\) can be solved for
\(\dot q^i=\dot q^i(q,p,t)\).
The Hamiltonian is
\[
  H(q,p,t)=p_i\dot q^i(q,p,t)-L(q,\dot q(q,p,t),t).
\]

\begin{figure}[htbp]
\centering
\begin{tikzpicture}[scale=1.0, >=Latex]
  \node[draw, rounded corners=2pt, minimum width=2.9cm, minimum height=0.9cm,
    fill=blue!5] (tq) at (0,0) {\(TQ:\ (q,\dot q)\)};
  \node[draw, rounded corners=2pt, minimum width=2.9cm, minimum height=0.9cm,
    fill=red!5] (tsq) at (5.0,0) {\(T^*Q:\ (q,p)\)};
  \draw[->, thick] (tq) -- node[above]
    {\(\mathbb F L:\ p_i=\partial L/\partial\dot q^i\)} (tsq);
  \draw[->, thick, bend left=18] (tsq.north west) to node[above]
    {regular inverse} (tq.north east);
  \node[align=center, text width=0.82\textwidth] at (2.5,-1.25)
    {The Legendre map preserves \(q\) but replaces velocity by the endpoint
    covector \(p\).\\It is a coordinate change only when the velocity Hessian is
    nonsingular.};
\end{tikzpicture}
\caption{The Legendre transform from tangent-bundle data to cotangent-bundle
data. Singular Lagrangians require constraints; the Hamiltonian construction in
this chapter assumes the regular case.}
\label{fig:legendre-transform}
\end{figure}

\begin{derivation}[Hamilton's equations from the Legendre transform]
Differentiate \(H\):
\[
\begin{aligned}
  \dd H
  &=
  \dot q^i\,\dd p_i
  +
  p_i\,\dd \dot q^i
  -
  \frac{\partial L}{\partial q^i}\dd q^i
  -
  \frac{\partial L}{\partial \dot q^i}\dd \dot q^i
  -
  \frac{\partial L}{\partial t}\dd t\\
  &=
  \dot q^i\,\dd p_i
  -
  \frac{\partial L}{\partial q^i}\dd q^i
  -
  \frac{\partial L}{\partial t}\dd t,
\end{aligned}
\]
because \(p_i=\partial L/\partial \dot q^i\). Thus
\[
  \frac{\partial H}{\partial p_i}=\dot q^i,
  \qquad
  \frac{\partial H}{\partial q^i}=-\frac{\partial L}{\partial q^i}.
\]
Using the Euler-Lagrange equation \(\dot p_i=\partial L/\partial q^i\), we get
\[
  \boxed{
    \dot q^i=\frac{\partial H}{\partial p_i},
    \qquad
    \dot p_i=-\frac{\partial H}{\partial q^i}
  }.
\]
\end{derivation}

The Legendre transform changes the bookkeeping of initial data. A Lagrangian
state \((q,\dot q)\) records a position and a velocity; a Hamiltonian state
\((q,p)\) records a position and the covector that measures how the action
changes under endpoint displacement. The symplectic form below is built from
that endpoint one-form, so Hamiltonian mechanics remembers the variational
origin of the momentum rather than merely renaming velocity.

\section{Phase Space As The Space Of Solutions}
\label{sec:phase-space-solutions}

A useful slogan is that, for a deterministic Hamiltonian system, phase space is
the same thing as the space of solutions.
More precisely, if \(H\) is smooth and the Hamiltonian vector field is complete
on the time interval under study, then every point
\[
  z_0=(q_0,p_0)\in T^*Q
\]
determines a unique trajectory \(z(t)=\Phi^t_H(z_0)\), and every trajectory is
obtained this way from its value at any one time. Thus phase space is the set
of all possible instantaneous states, but it is also a coordinate system on the
family of all classical histories. The caveat is that this identification uses
existence and uniqueness of the equations of motion; singular potentials,
collisions, constraints, or finite-time blowup can invalidate the global
statement even when the local Hamiltonian equations make sense.

\section{The Symplectic Form}
\label{sec:symplectic-form}

The cotangent bundle \(T^*Q\) carries a canonical one-form
\[
  \Theta=p_i\,\dd q^i.
\]
The canonical symplectic form is
\[
  \omega=-\dd\Theta=\dd q^i\wedge\dd p_i.
\]
In arbitrary local phase-space coordinates
\[
  x^m,\qquad m=1,\ldots,2n,
\]
one can write
\[
  \omega=\frac12\omega_{mn}(x)\,\dd x^m\wedge\dd x^n,
\]
where \(\omega_{mn}=-\omega_{nm}\), \(\det(\omega_{mn})\neq0\), and
\[
  \dd\omega=0.
\]
The closedness condition is what makes \(\omega\) locally equivalent to the
canonical form, while nondegeneracy is what lets the Hamiltonian determine a
unique vector field. In canonical coordinates \(x=(q,p)\),
\[
  (\omega_{mn})=
  \begin{pmatrix}
  0 & I_n\\
  -I_n & 0
  \end{pmatrix}.
\]
The Poisson bracket can be expressed intrinsically by the bivector inverse to
the symplectic form, with the sign convention fixed by Hamilton's equations
below. The point is that the bracket is not extra structure once \(\omega\) is
given.

The two defining properties of \(\omega\) have different jobs. Nondegeneracy
lets one solve \(\iota_{X_H}\omega=\dd H\) for a unique vector field, just as a
mass matrix lets one solve force equals mass times acceleration. Closedness is
the integrability condition behind canonical coordinates, Stokes-type
arguments, and the preservation of symplectic area. Hamiltonian mechanics needs
both.

The Hamiltonian vector field \(X_H\) is defined by
\[
  \iota_{X_H}\omega=\dd H.
\]
Here \(\iota_X\) denotes interior contraction: if \(\omega\) is a two-form, then
\((\iota_X\omega)(Y)=\omega(X,Y)\) for every vector field \(Y\).
If
\[
  X_H=A^i\frac{\partial}{\partial q^i}
  +
  B_i\frac{\partial}{\partial p_i},
\]
then
\[
  \iota_{X_H}\omega=A^i\,\dd p_i-B_i\,\dd q^i.
\]
Equating this with
\[
  \dd H=\frac{\partial H}{\partial q^i}\dd q^i
  +
  \frac{\partial H}{\partial p_i}\dd p_i
\]
gives
\[
  A^i=\frac{\partial H}{\partial p_i},
  \qquad
  B_i=-\frac{\partial H}{\partial q^i}.
\]
Thus \(X_H\) is Hamilton's equation written geometrically.

\begin{figure}[htbp]
\centering
\begin{tikzpicture}[scale=1.0, >=Latex]
  \draw[->] (-0.25,0) -- (6.9,0) node[right] {\(q\)};
  \draw[->] (3.35,-2.35) -- (3.35,2.35) node[above] {\(p\)};
  \node[below] at (0.25,0) {\(-\pi\)};
  \node[below] at (3.35,0) {\(0\)};
  \node[below] at (6.45,0) {\(\pi\)};
  \draw[blue!70!black, thick] (3.35,0) ellipse (0.85 and 0.55);
  \draw[blue!70!black, thick] (3.35,0) ellipse (1.45 and 0.95);
  \draw[blue!70!black, thick] (3.35,0) ellipse (2.05 and 1.35);
  \draw[red!70!black, thick, domain=-3.12:3.12, samples=180]
    plot ({\x+3.35},{1.62*cos(0.5*\x r)});
  \draw[red!70!black, thick, domain=-3.12:3.12, samples=180]
    plot ({\x+3.35},{-1.62*cos(0.5*\x r)});
  \draw[gray!65, thick, domain=0.45:6.25, samples=160]
    plot (\x,{1.9+0.18*cos(120*\x)});
  \draw[gray!65, thick, domain=0.45:6.25, samples=160]
    plot (\x,{-1.9+0.18*cos(120*\x)});
  \draw[->, thick] (4.15,0.78) -- (4.45,0.62) node[right] {\(X_H\)};
  \node[blue!70!black] at (1.45,1.75) {libration levels};
  \node[red!70!black] at (5.15,-1.9) {separatrix};
  \node[gray!65!black] at (1.15,-1.9) {rotation levels};
  \node at (3.35,-2.75) {Pendulum level sets: \(H(q,p)=\frac12p^2+1-\cos q\).};
\end{tikzpicture}
\caption{A schematic phase plane using the pendulum Hamiltonian. The
Hamiltonian vector field is tangent to energy level sets because
\(\dot H=\{H,H\}=0\).}
\label{fig:pendulum-phase-plane}
\end{figure}

\section{Poisson Brackets And Time Evolution}
\label{sec:poisson-time-evolution}

For \(f,g\in C^\infty(T^*Q)\), define
\[
  \Poisson{f}{g}
  =
  \frac{\partial f}{\partial q^i}
  \frac{\partial g}{\partial p_i}
  -
  \frac{\partial f}{\partial p_i}
  \frac{\partial g}{\partial q^i}.
\]
If \(z(t)=(q(t),p(t))\) follows Hamilton's equation, then
\[
  \frac{\dd}{\dd t}f(q(t),p(t),t)
  =
  \frac{\partial f}{\partial t}+\Poisson{f}{H}.
\]
A time-independent observable \(f\) is conserved exactly when
\[
  \Poisson{f}{H}=0.
\]

\begin{proposition}[Liouville volume preservation]
Hamiltonian flow preserves the phase-space volume form
\[
  \Omega=\frac{1}{n!}\omega^n.
\]
\end{proposition}

\begin{proof}
Cartan's formula gives
\[
  \Lie_{X_H}\omega=\dd(\iota_{X_H}\omega)+\iota_{X_H}\dd\omega.
\]
Since \(\iota_{X_H}\omega=\dd H\) and \(\dd\omega=0\),
\[
  \Lie_{X_H}\omega=\dd\dd H=0.
\]
Therefore \(\Lie_{X_H}\omega^n=0\), so the associated volume form is preserved.
\end{proof}

\section{Phase-Space Action}
\label{sec:phase-space-action}

Hamilton's equations can also be derived from the phase-space action
\[
  S_H[q,p]=\int_{t_0}^{t_1}\left(p_i\dot q^i-H(q,p,t)\right)\dd t.
\]
Varying \(q\) and \(p\) independently, with \(\delta q(t_0)=\delta q(t_1)=0\),
gives
\[
  \delta S_H
  =
  \int_{t_0}^{t_1}
  \left[
    \left(\dot q^i-\frac{\partial H}{\partial p_i}\right)\delta p_i
    -
    \left(\dot p_i+\frac{\partial H}{\partial q^i}\right)\delta q^i
  \right]\dd t.
\]
Since \(\delta q\) and \(\delta p\) are arbitrary in the interior, the
stationary curves are exactly Hamiltonian trajectories.

\section{Canonical Transformations}
\label{sec:canonical-transformations}

A diffeomorphism \(\Phi:T^*Q\to T^*Q\) is canonical, or symplectic, if
\[
  \Phi^*\omega=\omega.
\]
Canonical transformations preserve Poisson brackets:
\[
  \Poisson{f\circ\Phi}{g\circ\Phi}
  =
  \Poisson{f}{g}\circ\Phi.
\]
This is why canonical coordinates may change but Hamiltonian mechanics keeps
the same form. Later, action-angle variables will be canonical coordinates
adapted to integrable motion.

\begin{derivation}[Generating function for a canonical transformation]
Let \((q,p)\) be old canonical coordinates and \((Q,P)\) new ones. A convenient
type of generating function depends on the new momenta and old positions:
\[
  F=F(P,q).
\]
Impose
\[
  Q_i=\frac{\partial F}{\partial P_i}\bigg|_q,
  \qquad
  p_i=\frac{\partial F}{\partial q^i}\bigg|_P.
\]
Then
\[
  \dd F=Q_i\,\dd P_i+p_i\,\dd q^i.
\]
Taking another exterior derivative gives
\[
  0=\dd Q_i\wedge\dd P_i+\dd p_i\wedge\dd q^i.
\]
Therefore
\[
  \dd Q_i\wedge\dd P_i
  =
  \dd q^i\wedge\dd p_i.
\]
The canonical symplectic form is preserved, so the transformation is canonical.
In coordinates this is the same cancellation one sees by expanding
\[
  \dd P_i\wedge\dd Q_i
  =
  \dd P_i\wedge
  \left(
    \frac{\partial^2F}{\partial P_i\partial q^j}\dd q^j
    +
    \frac{\partial^2F}{\partial P_i\partial P_j}\dd P_j
  \right),
\]
where the pure \(\dd P_i\wedge\dd P_j\) term vanishes because second
derivatives are symmetric while the wedge product is antisymmetric.
\end{derivation}

\section{Numerical Lesson}
\label{sec:numerical-lesson}

The equations in this chapter are exact statements about smooth systems. A
numerical integrator creates a discrete map
\[
  z_{n+1}=\Psi_{\Delta t}(z_n).
\]
For long-time Hamiltonian problems, whether \(\Psi_{\Delta t}\) preserves a
symplectic form or a modified Hamiltonian can matter more than the local order
of the method. This is why the pendulum, standard map, and asteroid demos are
not merely programming exercises; they are experiments in preserving structure.

The practical lesson is diagnostic. A good mechanics simulation reports not
only a trajectory but also the quantities that the continuous problem says
should be preserved or nearly preserved. Energy drift, angular-momentum drift,
phase-space area, and event criteria are part of the model audit trail. Without
them, a numerical plot can look convincing while solving the wrong mechanical
problem.

\section{What To Take Forward}
\label{sec:lag-ham-take-forward}

The rest of the notes repeatedly reuse four lessons from this chapter.
\begin{itemize}
\item The configuration space \(Q\) is part of the model. Choosing bad
coordinates can obscure a simple system, but it cannot change the system.
\item Boundary terms are information, not clutter. They identify canonical
momenta, forces, and conserved quantities.
\item Symmetry produces conserved quantities through Noether's theorem, and
those quantities are often the coordinates of a reduced problem.
\item Hamiltonian mechanics is geometry on \(T^*Q\): the symplectic form
determines Hamilton's equations, volume preservation, Poisson brackets, and the
meaning of a canonical numerical method.
\end{itemize}
